%-----------------------------------------------------------------------
\subsection{Instance 7: Saliency Maps (GradCAM) Under Weight-Space
  Multiplicity}
\label{sec:instance-gradcam}
%-----------------------------------------------------------------------

Gradient-weighted Class Activation Mapping (GradCAM)
\citep{selvaraju2017gradcam} is the dominant saliency method for
convolutional neural networks: it highlights spatial regions that drive a
classification decision by weighting the final convolutional feature maps
by their class-specific gradients.  GradCAM saliency maps are widely used
for model debugging, clinical imaging, and regulatory documentation.  We
show that no saliency map explanation can be faithful, stable, and decisive
under weight-space underspecification.

\paragraph{The explanation system.}
Let $H \times W$ denote the spatial resolution of the final convolutional
feature maps.  We instantiate the abstract framework
(Definition~\ref{def:explanation-system}) as the 6-tuple
$\mathcal{S}_{\mathrm{sal}} = (\Params_{\mathrm{sal}},\,
\sH_{\mathrm{sal}},\,Y_{\mathrm{sal}},\,\mathsf{obs}_{\mathrm{sal}},\,
\mathsf{exp}_{\mathrm{sal}},\,\incomp_{\mathrm{sal}})$ where:
\begin{itemize}
  \item $\Params_{\mathrm{sal}}$ is the space of CNN weight configurations
    $\theta$ (convolutional filters, batch-norm parameters, classifier head);
  \item $\sH_{\mathrm{sal}} = [0,1]^{H \times W}$ is the space of
    ReLU-activated GradCAM heatmaps, normalized to $[0,1]$;
  \item $Y_{\mathrm{sal}}$ is the space of input-to-output prediction
    functions $f_\theta : \mathcal{X} \to \mathcal{Y}$;
  \item $\mathsf{obs}_{\mathrm{sal}}(\theta) = f_\theta$ records the
    observable prediction behavior;
  \item $\mathsf{exp}_{\mathrm{sal}}(\theta)$ is the GradCAM heatmap
    produced by $\theta$ on a fixed input image $x_0$; and
  \item $(\mathsf{h}_1, \mathsf{h}_2) \in \incomp_{\mathrm{sal}}$ when
    $\operatorname{argmax}_{(i,j)}\, \mathsf{h}_1 \neq
    \operatorname{argmax}_{(i,j)}\, \mathsf{h}_2$---the two heatmaps
    place peak activation at different spatial locations.
\end{itemize}

\paragraph{The Rashomon property.}
Sanity checks for saliency methods \citep{adebayo2018sanity} demonstrate
that saliency maps from networks that differ only in their random
initialization can diverge substantially while maintaining identical test
accuracy.  \citet{hooker2019benchmark} further show that saliency-guided
feature removal affects models inconsistently across retraining runs.
These findings establish the Rashomon property for
$\mathcal{S}_{\mathrm{sal}}$: there exist weight configurations $\theta_1,
\theta_2$ with $f_{\theta_1} \approx f_{\theta_2}$ yet peak GradCAM
activation at different spatial locations.

\begin{proposition}[Instance 7: Saliency Map Impossibility]
\label{prop:gradcam-impossibility}
If $\mathcal{S}_{\mathrm{sal}}$ has the Rashomon property, then no
saliency map explanation can simultaneously be faithful, stable, and
decisive.
\end{proposition}

\begin{proof}
Direct application of Theorem~\ref{thm:universal} to
$\mathcal{S}_{\mathrm{sal}}$.
\end{proof}

\paragraph{Empirical illustration.}
We validate Proposition~\ref{prop:gradcam-impossibility} by creating 10
perturbed ResNet-18 models on CIFAR-10 ($\sigma = 0.0005$ Gaussian weight
noise).  With \textbf{78.8\% prediction agreement} (confirming approximate
functional equivalence; prediction agreement is 78.8\%, reflecting
sensitivity of ImageNet-pretrained ResNet-18 to perturbation on CIFAR-10),
the \textbf{peak flip rate} is \textbf{9.6\%}
[95\% CI: 8.8--10.5\%]: nearly one in ten model pairs places peak GradCAM
activation at a different spatial location despite the same prediction.
The mean pairwise spatial IoU (top-20\% threshold) is \textbf{95.4\%}
[CI: 95.3--95.5\%], confirming high but imperfect overlap.
The negative control (identical weights) yields 100\% IoU and 0\% flip
rate.  The resolution test (averaged GradCAM heatmaps) yields a mean IoU
of \textbf{96.9\%}---a \textbf{+1.5\%} improvement over pairwise IoU,
confirming that orbit averaging reduces instability.
Full results appear in Table~\ref{tab:gradcam}.

\begin{table}[ht]
\centering
\caption{GradCAM Saliency Map Instability (ResNet-18, CIFAR-10).
  Ten functionally equivalent models (pretrained ResNet-18 + Gaussian perturbations, $\sigma=0.02$).
  Spatial IoU uses top-20\% threshold; peak flip rate counts pairs where argmax differs by $>$2 pixels.
  Negative control uses identical weights (model 0 for all 10 slots).
  Resolution test averages heatmaps; mean IoU vs.\ average exceeds pairwise IoU.
  95\% bootstrap CIs in parentheses.}
\label{tab:gradcam}
\begin{tabular}{lccc}
\toprule
Metric & Positive & Control & Resolution \\
\midrule
Prediction agreement & 0.788 (0.752, 0.822) & 1.000 & --- \\
Mean pairwise IoU & 0.954 (0.953, 0.955) & 1.000 (1.000, 1.000) & 0.969 (0.968, 0.971) \\
Peak flip rate & 0.096 (0.088, 0.105) & 0.000 (0.000, 0.000) & --- \\
\bottomrule
\end{tabular}
\end{table}


\paragraph{Moderate instability compared to attention.}
The 9.6\% peak flip rate for GradCAM is substantially lower than the 60\%
argmax flip rate for attention maps (Section~\ref{sec:instance-attention}).
This disparity reflects the structural constraint imposed by spatial
locality: GradCAM operates on $7 \times 7$ or $8 \times 8$ feature maps
where neighboring activations are highly correlated, whereas token-level
attention distributes mass across an unstructured sequence.  Spatial
smoothness constrains the degree to which perturbations can rearrange
saliency, producing a tighter Rashomon set for the peak-location
incompatibility relation.  Nevertheless, the impossibility still holds:
even moderate instability suffices for the Rashomon property, and the 9.6\%
flip rate is statistically significant with a bootstrap CI excluding zero.
